\section{The Per-Cell State of the Art}
\label{sec:sota}

In this section, we describe how existing algorithms compute one cell of the spectrum, and why repeating them across a range is expensive.
%
Understanding the per-cell cost is what motivates the shared build and the transfer theory of the following sections.

\stitle{Peeling one cell.}
The standard algorithm for a single cell $(r,s)$ follows Definition~\ref{def:core} directly~\cite{firstNucleus, nuclearTWEB}.
%
It enumerates every $r$-clique, computes the initial support of each by counting the $s$-cliques that contain it, and then repeatedly removes an $r$-clique of minimum support, decrementing the supports of the $r$-cliques that share an $s$-clique with the removed one, until the graph is empty.
%
The core value of an $r$-clique is the support at which it is removed.
%
The cost is dominated by two quantities: the number of $r$-cliques, which sets the size of the working state, and the number of $s$-cliques, which sets the counting cost.

\stitle{Counting without enumeration.}
The state-of-the-art algorithm \cnd~\cite{LocalNuclear} improves the counting step.
%
Instead of enumerating the $s$-cliques around each $r$-clique, it maintains a counting structure that returns the support of an $r$-clique by combinatorial aggregation over a clique index, avoiding explicit $s$-clique enumeration inside the cell.
%
This makes \cnd{} the fastest per-cell algorithm on dense graphs, and its per-cell counting cost is nearly flat in $s$: on \daAST{}, the cells $(4,5)$ through $(4,8)$ each take about $40$ seconds.

\stitle{The per-cell wall.}
Both frameworks share the same structural cost: they hold state for every $r$-clique of the cell.
%
The number of $r$-cliques grows sharply with $r$ and with graph density, so the working state grows with it.
%
On \daDBL{} at $r{=}5$, one cell holds enough state to reach $257$GB of memory, and on \daWIT{} every cell exceeds a $300$GB budget.
%
Because nothing is shared between cells, a spectrum of $w$ cells pays this cost $w$ times: on \daDBL{} at $r{=}5$, the four cells $s{=}6,\dots,9$ cost $4{,}261$ seconds in total, and on \daHEP{} the single cell $(3,7)$ alone exceeds a two-hour budget.

\begin{remark}
The per-cell wall is not an implementation detail that a faster counting structure removes.
%
It follows from computing each cell from scratch: the working state is proportional to the $r$-clique count of the cell, and the count is paid once per cell.
%
The rest of this paper removes the wall by computing cells from one another, so that only a boundary cell and a small residue are ever peeled.
\end{remark}
